\documentclass[10pt]{article}
\usepackage[margin=0.85in]{geometry}
\usepackage[T1]{fontenc}
\usepackage{lmodern}
\usepackage{amsmath,amssymb,booktabs,tabularx,array}
\usepackage{algorithm,algorithmic}
\usepackage[hidelinks]{hyperref}
\usepackage{microtype}
\newcommand{\relu}[1]{\left[#1\right]_+}
\newcommand{\norm}[1]{\left\lVert#1\right\rVert}
\newcommand{\rms}{\operatorname{RMS}}
\newcommand{\mean}{\operatorname{mean}}
\newcommand{\clip}{\operatorname{clip}_{[0,1]}}
\newcommand{\pathcode}[1]{\nolinkurl{#1}}
\newcolumntype{Y}{>{\raggedright\arraybackslash}X}
\setlength{\parindent}{0pt}
\setlength{\parskip}{5pt}
\title{Current Guidance: Algorithm, Objective, and Generality}
\author{Implementation note for ParticleDiffuser}
\date{September 15, 2026}
\begin{document}
\begin{center}
{\Large\bfseries Current Guidance: Algorithm, Objective, and Generality}\\[4pt]
ParticleDiffuser implementation note \quad September 15, 2026
\end{center}

\textbf{Scope.} This describes the \pathcode{noise_opt} and
\pathcode{select_refine} paths used in the September 15 guidance campaign,
checked against code revision \pathcode{b484a309611f}.
It is an implementation description, not a claim that the resulting plans are
physically valid. The campaign uses a frozen, unconditional, bidirectional,
pose-native UNet. The current unconditional entrypoint supports one entirely
rigid scene at a time; it is not presently a general cloth/MPM planner.

\textbf{Main distinction.} The optimizer is reusable: differentiate a trajectory
cost through a frozen generator and update its initial noise. The current
pick/place objective is specialized: it contains eight auxiliary picking
penalties, or eleven for placing, in addition to task goals and noise
regularization. Combining those penalties into one symbol does not make them
a general physical law.

\section{Core noise optimization}
Let $c$ be clean observed history plus scene metadata, $g$ a requested goal,
and $G_\theta$ the full windowed sampler. Its output $\tau$ includes both object
and gripper motion. Model weights $\theta$ are fixed. For each rollout window
$w$, save particle noise $z_w^p$, pose noise $z_w^q$, and model-routing RNG state
$r_w$. Only the pose noise is changed:
\begin{align}
 \tau(\Delta)&=G_\theta\bigl(c,\{z_w^p,z_w^q+M_w\odot\Delta_w,r_w\}_{w=1}^W\bigr),\\
 J(\Delta)&=L_{\rm task+aux}(\tau(\Delta);g)
       +\lambda_z\frac{\sum_w\norm{M_w\odot\Delta_w}_F^2}{\sum_w n_w},
 \qquad
 \rms_{M_w}(\Delta_w)\leq\rho\quad\text{for each }w.
\end{align}
Here $M_w$ selects valid pose channels and $n_w$ counts its active scalar entries,
including time. The default scope includes passive object and actuator channels;
the actuator-only ablation masks out passive updates. The particle-noise arrays
and routing RNG are held fixed. The regularizer penalizes change from the
initial draw; it is \emph{not} a learned feasibility score or an exact Gaussian
log-prior on the resulting noise.

\begin{algorithm}[H]
\caption{Current complete-rollout initial-noise optimization}
\begin{algorithmic}[1]
\REQUIRE Frozen generator $G_\theta$; history $c$; goal $g$; updates $I$;
step size $\eta$; cap $\rho$; change weight $\lambda_z$
\STATE Use a supplied initial rollout, or draw an ordinary one; retain $(z_w^p,z_w^q,r_w)$
\STATE Initialize $\Delta_w\gets0$ for all windows; initialize Adam; $J^*\gets+\infty$
\FOR{$i=0,\ldots,I$}
  \STATE Regenerate \emph{all} windows from $c$, using the saved noise plus $M\odot\Delta$
  \STATE Replay $r_w$ in each window; keep gradients through carried predicted history
  \STATE Compute $J(\Delta)$ on the complete decoded trajectory
  \IF{$J(\Delta)<J^*$}
    \STATE $J^*\gets J(\Delta)$; save a detached copy of this rollout and its noise
  \ENDIF
  \IF{$i<I$}
    \STATE Backpropagate through every window and denoising step to $\{\Delta_w\}$
    \STATE $\{\Delta_w\}\gets\operatorname{AdamStep}(\{\Delta_w\},\nabla_\Delta J,\eta)$
    \FOR{each window $w$}
      \STATE $\Delta_w\gets M_w\odot\Delta_w$
      \STATE $\Delta_w\gets\Delta_w\min\{1,\rho/\max(\rms_{M_w}(\Delta_w),10^{-12})\}$
    \ENDFOR
  \ENDIF
\ENDFOR
\STATE Run the separate geometric audit of the best saved trajectory
\RETURN Best saved trajectory, optimization trace, and geometric audit
\end{algorithmic}
\end{algorithm}

\clearpage
\section{What one differentiable rollout does}
The multiwindow procedure is one optimization over the whole requested future.
It does not execute an action or acquire a new observation between windows.
The next window conditions on the previous window's generated history.

\begin{algorithm}[H]
\caption{Windowed sampler $G_\theta$ inside each objective evaluation}
\begin{algorithmic}[1]
\REQUIRE Observed context $c$; horizon $H$; window length $F$; noise and RNG per window
\STATE Initialize accumulated particles and rigid poses from the clean observations
\WHILE{fewer than $H$ future frames have been produced}
  \STATE Build the next window from the last $C$ accumulated context frames
  \STATE Shorten its future to the remaining horizon; hide all future action conditioning
  \STATE Initialize particle/pose diffusion states with that window's saved noise
  \STATE Apply articulation masks and restore the window's saved routing RNG
  \FOR{successive rows $k_m,k_{m+1}$ of the diffusion-forcing schedule}
    \STATE Skip a row pair only if every frame's timestep is unchanged
    \STATE Evaluate the frozen model with clean context and current noisy future states
    \STATE Convert model predictions to velocities using the trained output convention
    \STATE Update each future frame by Euler integration toward the next noise level
    \STATE Reapply the valid pose-channel mask
  \ENDFOR
  \STATE Decode normalized deltas into physical poses and rigid-template particles
  \STATE Append the window's new frames; preserve its computational graph
\ENDWHILE
\RETURN Accumulated object and actuator trajectories
\end{algorithmic}
\end{algorithm}

For either diffusion state $x$ (particles or poses), the Euler update is
\[
 x_f^{m+1}=x_f^m-\frac{k_{m,f}-k_{m+1,f}}{T_{\rm train}}
                  v_\theta(x^m,c,k_m)_f.
\]
This is the ordinary sampler inside an outer optimization. The campaign does
not add a cost gradient at every denoising step; that is a separate optional
method in the repository. Freezing routing RNG stabilizes random representative
point choices; it does not freeze coordinate-dependent spatial neighborhoods.

\textbf{Concrete settings.} Context $C=4$, window length $F=16$, nominal
denoising steps $25$, pyramid schedule, and physical timestep $0.05$ s.
$H=36$ uses windows $16+16+4$ (1.8 s); $H=68$ uses $16+16+16+16+4$ (3.4 s).
The pyramid schedule can require more model calls than its nominal step count.
The main runs use $I=64$, $\eta=0.03$, $\rho=0.5$, $\lambda_z=0.001$, full-frame
updates, constant learning rate, no gradient normalization, and gradients
across window boundaries. Smaller-budget/cap ablations are recorded separately.

\textbf{Selection and optional variants.} The main campaign has geometric
screening and augmented-Lagrangian constraints \emph{off}: the minimum weighted
objective, including the change penalty, selects the iterate. Geometry is
audited afterward and can fail. An optional screen elsewhere ranks full
geometric acceptance first, non-goal geometric acceptance second, then
interaction violation and loss. That screen is still not physics execution.
For \pathcode{select_refine}, first sample $K=16$ complete trajectories, select
the minimum task-plus-auxiliary cost, and run Algorithm 1 for 24 updates around
that selected draw. No fixed grasp proposal or prescribed path is supplied.

\textbf{Cost of the procedure.} The bottle/teacup campaign measured about
23.4 minutes for 36 frames/64 updates and 43.2 minutes for 68 frames/64 updates
on one PRO 6000 in FP32 with activation checkpointing. The former uses about
14,040 model calls including recomputation. These are measured research-run
costs, excluding loading/final audits/rendering, not optimized online latency.

\clearpage
\section{What is actually in the objective?}
For the mesh-distance path used in this campaign,
\begin{align}
 L_{\rm task+aux}
 &=w_g L_g+\mathbb{1}_{\rm place}\,w_R L_R
       +\alpha\sum_{j\in\mathcal A}w_jL_j
       +\mathbb{1}_{\rm obstacles}\,w_oL_o,\\
 \mathcal A_{\rm pick}
 &=\{\text{contact, rise, slip, penetration, object floor, gripper floor,
                   joint limits, joint speed}\},\nonumber\\
 \mathcal A_{\rm place}
 &=\mathcal A_{\rm pick}\cup\{\text{rest, release, settle}\}.\nonumber
\end{align}
The default physical scale is $\alpha=1$. The zero-physical-cost ablation sets
$\alpha=0$ while retaining position/rotation goals, noise regularization and any
obstacle objective. Picking and placing also use different contact formulas.

\small
\begin{tabularx}{\linewidth}{@{}l r Y@{}}
\toprule
Term & Default weight & Meaning and assumptions\\
\midrule
Position goal & 1 & Mean or maximum distance over the last six frames; uses object origin.\\
Orientation goal (place) & 0.04 & Mean geodesic rotation error; coefficient is metres/radian.\\
Finger contact & 2 & Encourage both fingers near the lifted object; pick also rewards final proximity.\\
Unsupported rise & 5 & Penalize upward object displacement when fingers were far at the previous frame.\\
Slip & 1 & Penalize change in object origin in the hand frame while close and lifted.\\
Object--gripper penetration & 5 & Mean and peak excess convex-component overlap.\\
Object below support & 5 & Mean and peak intrusion below a supplied horizontal plane.\\
Gripper below support & 5 & Same plane check for hand and fingers.\\
Finger joint limits & 10 & Mean and peak excess beyond calibrated prismatic travel.\\
Finger joint speed & 10 & Mean and peak excess per-frame travel beyond calibrated speed times $\Delta t$.\\
Rest (place) & 5 & Put the object's lowest mesh vertex on the support plane.\\
Release (place) & 2 & Separate both fingers from the object in the final interval.\\
Settle (place) & 2 & Reduce final object-origin translation between frames.\\
Obstacle clearance (optional) & 5 & Mean and peak shortfall from a 5 mm supplied-box clearance.\\
Noise-change regularization & 0.001 & Mean squared active pose-noise change, plus a separate hard RMS cap.\\
\bottomrule
\end{tabularx}
\normalsize

\textbf{Counting honestly.} Picking has one goal, eight auxiliary costs and
one noise penalty: ten scalar contributions. Placing has two goals, eleven
auxiliary costs and one noise penalty: fourteen. Obstacles add one more. There
are additional tolerances, gates and hold lengths; the weights are not the
only hand-chosen quantities. The table lists the base configuration, not the
stronger-weight ablations. The slip quantity includes changing hand coordinates
but has no explicit object-relative-rotation term; settling only measures
translation. The code's diagnostic key \pathcode{hand_step} refers to this
object-in-hand coordinate change, not simply hand speed.

\textbf{Why the generality concern is valid.} Noise optimization itself is a
general numerical mechanism. Joint travel and collision geometry can be
reused when a new robot/environment supplies those constraints. However,
``upward without two nearby fingers'', ``stay fixed relative to the hand'',
and ``rest on a horizontal plane then release'' encode particular manipulation
strategies. They are inappropriate for many pushes, throws, sliding contacts,
cloth interactions or granular flows. No grasp proposal is required, but the
objective still contains a two-finger grasp assumption and a vertical support
assumption. These should be disclosed as task-specific regularizers.

\clearpage
\appendix
\section{Exact mesh-path cost definitions}
All lengths below are in metres; write $[x]_+=\max(x,0)$ and
$A(v)=\mean(v)+\max(v)$ over the specified entries. Let $t=0$ denote the last
observed anchor and $t=1,\ldots,H$ future states. Let $p_t,R_t$ be the object
origin and rotation, $h_t,Q_t$ the hand origin and rotation, and $b_t$ the
object's lowest transformed collision-mesh vertex. Let $b^k_t$ be the lowest
vertex of gripper link $k$, $z_s$ the support plane, and
$\mathcal T=\{H-5,\ldots,H\}$ the six-frame goal interval.

Let $s_{t,k}$ be the signed object-to-gripper distance for hand or finger $k$
(negative for overlap). Distances use the minimum across supplied convex
collision components. Define the farther-finger gap, lift gate, missing-contact
gate, and object displacement in hand coordinates:
\begin{align}
 d_t&=\max_{k\in\{f_1,f_2\}}[s_{t,k}]_+,
 &e_t&=[d_t-0.002]_+,\\
 a_t&=\clip\!\left(\frac{b_t-z_s-0.003}{0.02}\right),
 &m_t&=\clip\!\left(\frac{d_{t-1}-0.002}{0.01}\right),\\
 u_t&=Q_t^\top(p_t-h_t),
 &v_t&=\norm{u_t-u_{t-1}}_2.
\end{align}
The 2 mm contact tolerance, 3 mm support tolerance, 20 mm lift ramp and 10 mm
missing-contact ramp are explicit design choices. The current pick contact
cost has a four-frame terminal proximity term, even though the goal uses six.

\subsection{Task goals and grasp-related penalties}
For the main point-target experiments,
\begin{align}
 L_g&=\operatorname{reduce}_{t\in\mathcal T}\norm{p_t-g}_2,
 &&\operatorname{reduce}\in\{\mean,\max\},\\
 L_R&=\mean_{t\in\mathcal T}\operatorname{angle}(R_g^\top R_t),
 &&\text{placing only},\\
 L_{\rm contact}^{\rm pick}
 &=\mean_{t=1:H}(a_te_t)+\mean_{t=H-3:H}e_t,\\
 L_{\rm contact}^{\rm place}&=\mean_{t=1:H}(a_te_t),\\
 L_{\rm rise}&=\mean_{t=1:H}\bigl([p_{t,z}-p_{t-1,z}]_+m_t\bigr),\\
 L_{\rm slip}&=\mean_{t=1:H}\bigl(v_t(1-m_t)a_t\bigr).
\end{align}
Early comparisons use mean goal error; the selected jar/bowl replication uses
maximum held-position error. The API also supports position regions for
picking; that replaces the point distance with the supplied region's distance.
None of these proximity gates tests friction, contact force or force balance.

\subsection{Collision and support penalties}
With $P_{t,k}=[-s_{t,k}-0.002]_+$ for all three gripper links,
\begin{align}
 L_{\rm penetration}&=\mean_{t=1:H}\max_k P_{t,k}
                       +\max_{t=1:H,k}P_{t,k},\\
 L_{\rm floor}^{\rm object}&=A\bigl([z_s-b_t-0.003]_+:t=1,\ldots,H\bigr),\\
 L_{\rm floor}^{\rm gripper}&=A\bigl([z_s-b^k_t-0.003]_+:t=1,\ldots,H;
                                  k\in\{h,f_1,f_2\}\bigr).
\end{align}
FCL supplies distances and witness points. The implementation preserves the
distance value and uses frozen witnesses for a local first derivative with
respect to poses. Collision-component selection is not differentiated, and
feature/contact switches are nonsmooth. This is a geometric derivative, not
differentiation through a contact simulator. Overlapping convex components
also need not define the exact penetration distance of their union.

\clearpage
\section{Joint, placement and obstacle details}
For finger $j$, the calibrated prismatic coordinate is reconstructed from its
predicted position relative to the hand:
\[
 q_{t,j}=q_{0,j}+\left\langle
 Q_t^\top(p^{f_j}_t-h_t)-\ell_{0,j},\,a_j^{\rm joint}\right\rangle.
\]
Here $\ell_{0,j}$ is its anchor offset in hand coordinates and $a_j^{\rm joint}$
the calibrated axis. With travel limits $[q_j^-,q_j^+]$ and speed limit
$\dot q_j^{\max}$,
\begin{align}
 L_{\rm joint}&=A\bigl([\max(q_j^- -q_{t,j},q_{t,j}-q_j^+)]_+:t=1:H,j\bigr),\\
 L_{\rm joint\,speed}&=A\bigl([|q_{t,j}-q_{t-1,j}|-\dot q_j^{\max}\Delta t]_+
                                      :t=1:H,j\bigr).
\end{align}
The latter is a length-valued excess per transition, not a squared velocity
cost. There is no corresponding generic root-hand speed/acceleration penalty
in this objective.

The additional placement terms are
\begin{align}
 L_{\rm rest}&=\mean_{t\in\mathcal T}|b_t-z_s|,\\
 L_{\rm release}&=\mean_{t\in\mathcal T}
                       [0.005-\min_{j\in\{f_1,f_2\}}s_{t,j}]_+,\\
 L_{\rm settle}&=\mean_{t=H-4:H}\norm{p_t-p_{t-1}}_2.
\end{align}
Release uses signed distance, so overlap is penalized as well as insufficient
separation. The five settling transitions connect the last six states.

For optional supplied static boxes, let $d^{o}_{\ell,k,b}$ be signed clearance
of moving body $k$ from box $b$ at interpolation probe $\ell$. Include every
supplied moving body and interpolate translation and shortest rotation from
the observed anchor through the future. Excluding the fixed anchor from the
optimization cost,
\[
 L_o=A\bigl([0.005-d^{o}_{\ell,k,b}]_+:\ell>0,k,b\bigr).
\]
The campaign uses three probes per transition for optimization and ten for
the final geometric audit, which includes the anchor. This is not exact
continuous collision detection. The obstacle boxes enter the objective and
replay world; the frozen network does not receive them as conditioning.

\section{What the optimizer does not establish}
The default noise mask allows the optimizer to change the passive object's
future as well as the hand's future. A small joint-trajectory cost does not
establish that executing the hand causes that object motion. Restricting
updates to actuator noise leaves the model jointly generative; it does not
convert it into a verified conditional dynamics model.

The geometric audit, the optimization objective and independent execution are
distinct. The MuJoCo diagnostic commands a floating hand and articulated
fingers while the object evolves dynamically. Converting generated achieved
finger coordinates into motor commands uses a separate calibrated preload
policy. The optional latched policy is a controller ablation, not a loss term.
Neither that diagnostic nor this objective establishes source-simulator or
hardware validity.

The campaign's new jar/bowl goals give three model position hits out of four,
but no complete independent-execution passes. These are previously inspected
fixtures and exploratory goals, not a population success estimate. The
zero-physical-cost comparison also produced unsupported predicted lifting and
increased penetration. Simply deleting the auxiliary terms is therefore an
honest baseline to test, not an established fix.

\clearpage
\section{A simpler formulation to investigate, not the current method}
The appropriate general claim is \emph{task-cost optimization through a
generative trajectory model}. The current objective should be described as a
pick/place instantiation with explicit geometric regularization. It should not
be described as a universal physical-feasibility objective.

A useful separation for future work is:
\begin{enumerate}
 \item \textbf{Task request:} target position/orientation or another measurable
       outcome. The task must specify success, but need not specify a grasp
       phase schedule or a motion trajectory.
 \item \textbf{Declared robot/environment feasibility:} joint travel, speed and
       collision constraints supplied by the robot and environment. These
       still require geometry, calibration and tolerances.
 \item \textbf{Motion consistency:} whether the requested actuator motion
       produces compatible passive-object motion across independently sampled
       futures and, separately, independent execution.
\end{enumerate}

For example, if a qualified conditional sampler $x^{(s)}\sim
p_\theta(x\mid c,u)$ were available for a shared actuator trajectory $u$, an
experimental objective could take the form
\[
 \min_u\ \frac{1}{S}\sum_{s=1}^{S}L_g(x^{(s)})+\lambda_u R(u)
 \quad\text{subject to declared robot/environment constraints.}
\]
This would reduce dependence on hand-designed ``rise'', ``slip'', and
``release'' surrogates only if the conditional dynamics and action interface
are reliable. The existing actuator-inpainting diagnostic is not a qualified
conditional sampler, and this multi-draw objective is not implemented by the
current campaign. Consistent wrong predictions would still pass an internal
consistency check; execution validation remains necessary.

The immediate comparison should expose, rather than hide, the design choices:
goal plus noise regularization; add robot/environment geometry; then add the
task-specific interaction terms. Compare equal planning horizons and measured
model-evaluation budgets, retain independent evaluation seeds, and report
goal achievement separately from geometric and execution failures. No
generality or improvement claim follows from merely reducing the term count.

\section{Source map and reproducibility}
All source files below are in \pathcode{guidance/}, at the revision on page 1.

\small
\begin{tabularx}{\linewidth}{@{}Y Y@{}}
\toprule
Source & Responsibility\\
\midrule
\pathcode{unconditional.py} & \pathcode{noise_opt} (189),
 \pathcode{sample_rollout} (65), \pathcode{select_then_refine} (359).\\
\pathcode{sampler.py}, \pathcode{run_unconditional.py} & Euler sampler,
 frozen-model CLI and objective/audit wiring.\\
\pathcode{pick_objective.py}, \pathcode{physical_objective.py} & Base costs,
 role/geometry features, held goals, joints and support.\\
\pathcode{mesh_physical.py}, \pathcode{place_objective.py} & Mesh replacements
 and placement-specific terms.\\
\pathcode{obstacles.py}, \pathcode{collision_distance.py} & Box costs/audits
 and FCL witness-based derivatives.\\
\pathcode{pick_validation.py}, \pathcode{replay_control.py} & Separate
 geometric selection rules and replay command policies.\\
\bottomrule
\end{tabularx}
\normalsize

Campaign records:\\
\pathcode{validation/guidance_push_2026-09-15/}\\
\pathcode{COMPUTE.md}, \pathcode{REPLICATION.md},
\pathcode{PHYSICAL_ABLATION.md}, and individual \pathcode{protocol.json} files.
The checkpoint SHA256 begins \pathcode{99ec6cc6102a0d65}.
Earlier immutable campaign snapshots predate some optional CLI switches;
their default behavior is represented above. Staged grasp proposals,
augmented-Lagrangian constraints, local denoising guidance, beam search and
feedback replanning are separate optional paths, not silently included here.

\textbf{Build:} Run \pathcode{pdflatex -interaction=nonstopmode -halt-on-error}
on this file (twice for references). It has no external figures or bibliography.
\end{document}
